\scanpage{41}

\(f:X\longrightarrow Y\) morphisme propre

\diagram{\begin{tikzcd}[row sep=large,column sep=large]
K(X)\arrow[r,"f_!"]\arrow[d,"\widetilde C_X"']&K(Y)\arrow[d,"\widetilde C_Y"]\\
\widetilde A(X)\arrow[d,"h_X"']&\widetilde A(Y)\arrow[d,"h_Y"]\\
A(X)\otimes\mathbb Q\arrow[d,"\tau_h(X)\mathbin{\cdot}"']&A(Y)\otimes\mathbb Q\arrow[d,"\tau_h(Y)\mathbin{\cdot}"]\\
A(X)\otimes\mathbb Q\arrow[r,"f_*"']&A(Y)\otimes\mathbb Q
\end{tikzcd}}
\[
\boxed{\tau_h(Y)\varphi_Y(f_!(\xi))=f_*(\tau_h(X)\varphi_X(\xi))}\tag{1}
\]

\uline{Exemples}

1) \(f\) une injection. \uline{Il y a alors une formule meilleure}, donnant
\[
\boxed{C^j_Y(f_!(\xi))=f_*\left[\frac{C^{q+j}_{\deleted{\unsure{X}}}(\deleted{\unsure{\xi}}\widetilde C(\xi)*\lambda_{-1}(\widetilde C(T_{Y/X})))}{(-1)^q C^q_Y(T_{Y/X})}\right]}\tag{2}
\]
\deleted{\unsure{et}} \unsure{et}

et (1) prend la forme
\[
\boxed{\varphi_Y(f_!(\xi))=f_*(\varphi_X(\xi)\tau_h(T_{Y/X})^{-1})}\tag{1bis}
\]

Noter qu'ici \(f_!(\unsure{\gamma_X(F)|\sigma_X(F)})=\gamma_Y(F)\).

2) \(f\) est une application fibrée. Supposons \(E\) fibré vectoriel sur \(X\), alors on a
\[
f_!(\gamma_X(E))=\sum(-1)^i\cdot\gamma_Y(E^{(i)})
\]
où \(E^{(i)}\) est le fibré vectoriel sur \(Y\) dont la fibre est le \(i\)ième groupe de cohomologie de la fibre de \(X\to Y\).

On a donc \(\varphi_Y(f_!(\gamma_X(E)))=\sum(-1)^i\varphi_Y(E^{(i)})\). Dans le cas où \(E^{(i)}=0\) pour \(i>0\) (p.ex. il en est ainsi si on remplace \(E\) par \(E(n)\) \inserted{\(n\) grand}, provenant d'une immersion de \(X\) comme loc. fermé dans un espace projectif) on détermine ainsi \(\varphi_Y(E^{(0)})\) donc \(\widetilde C_Y(E^{(0)})\) \unsure{mod. la torsion}\deleted{\(\unsure{\xi}\)}. \deleted{\unsure{Si on}} \unsure{On a} dans la formule

\margin{R}
\figure{\begin{tikzpicture}[scale=.2]\draw (0,0) circle (.35);\foreach \a in {0,45,...,315}{\draw (\a:.7) circle (.35);}\end{tikzpicture}}
\margin{ARGENTINA LINEN TYPEWRITER}
\figure{\begin{tikzpicture}[scale=.2]\draw (0,0) circle (.35);\foreach \a in {0,45,...,315}{\draw (\a:.7) circle (.35);}\end{tikzpicture}}
\margin{\unsure{A}}
\margin{40}
